\documentclass[12pt]{article}
\usepackage[margin=1in]{geometry}
\usepackage{microtype}
\usepackage{xurl}
\usepackage[table]{xcolor}
\usepackage{hyperref}
\hypersetup{
  colorlinks=true,
  urlcolor=blue,
  linkcolor=blue,
  citecolor=blue
}
\setlength{\parindent}{0pt}
\setlength{\parskip}{0.8em}
\setlength{\emergencystretch}{4em}
\sloppy

\begin{document}
\begin{center}
{\LARGE San Jose State University}\\[0.5em]
{\LARGE CS 147 - Computer Architecture}\\[0.5em]
{\LARGE Class Project - Phase 3}
\end{center}

\vspace{2em}

\subsection*{Overall Timeline and Grade Distribution}
\begin{center}
\begin{tabular}{|l|l|l|}
\hline
\textbf{Due Date} & \multicolumn{2}{l|}{\textbf{Project Component}} \\
\hline
2/26/26 & Design Review & (5\% of project grade) \\
\hline
3/10/26 & Phase \#1 & (15\% of project grade) \\
\hline
4/7/26 & Phase \#2 & (30\% of project grade) \\
\hline
4/16/26 & Phase \#2.1 & (10\% of project grade) \\
\hline
N/A & Phase \#2.2 & (Optional) \\
\hline
4/23/26 & Phase \#2.3 Caching & (10\% of project grade) \\
\hline
\rowcolor{black!8}
5/7/26 & Phase \#3 & (30\% of project grade) \\
\hline
\end{tabular}
\end{center}

\section*{Phase \#3 --- Final Processor Integration (30\% of project grade)}

\subsection*{Overview}

Phase 3 is the final stage of the processor project. In this phase you will integrate your
fully pipelined processor with a realistic memory system and implement several required
performance optimizations.

By this point, your design should already include:

\begin{itemize}
\item A working five-stage pipelined processor
\item A functioning cache system
\item Correct handling of pipeline hazards
\end{itemize}

Phase 3 focuses on integrating the full memory hierarchy and implementing the
forwarding and prediction mechanisms needed for efficient execution.

\subsection*{Required Features}

Your final processor must implement the following functionality:

\begin{itemize}
\item Two-way set-associative caches backed by multi-cycle memory
\item Register file bypassing
\item Forwarding from the beginning of the MEM stage to the beginning of the EX stage
\item Forwarding from the beginning of the WB stage to the beginning of the EX stage
\item Static branch prediction (predict not taken)
\end{itemize}

When implementing forwarding, ensure that your design does not create a combinational
path that effectively performs the work of multiple pipeline stages in a single cycle.
Violating the intended pipeline structure may result in loss of credit.

The IPC of your processor will also account for a small portion of this demo grade, so you
should attempt to eliminate unnecessary stall cycles where possible.

\subsection*{Assembling / Simulating}

For the WISC-SJ26 ISA specification, refer to:

\begin{quote}
\texttt{./assignments/project/common/writeup/wisc-sj26\_spec.pdf}
\end{quote}

An assembler for the WISC-SJ26 ISA is provided for your use. Sample test
programs are also provided, although you are strongly encouraged to write custom
tests to augment these. Be aware that some provided test programs were written
for a slightly different ISA specification and therefore may not work exactly as
advertised. If unexpected behavior occurs, you must determine whether the issue
is in your design or in the assumptions made by the test itself.

The expected workflow for debugging WISC-SJ26 assembly tests is:

\begin{enumerate}
\item Write a test in WISC-SJ26 assembly language.
\item Assemble it using:
\begin{verbatim}
./run assemble <program>.asm <program>
\end{verbatim}
\item Simulate the generated image using:
\begin{verbatim}
./run simulate <program>_all.img
\end{verbatim}
\item Examine the simulator trace to confirm that each instruction and register
update matches your expectations.
\end{enumerate}

The simulator prints an instruction-by-instruction execution trace, which should
be used to confirm that your program is doing exactly what you intended before
you rely on it as evidence that the hardware is correct.

\subsection*{Cache Files To Carry Forward}

Your final cache implementation work for this project path should live in:
\begin{quote}
\texttt{./demo2/verilog/phase2\_3/cache\_assoc/}
\end{quote}

For Phase 3, the carried-forward cache should be the \textbf{two-way set-associative}
version that serves as the final cache design target.

The carried-forward cache files should include:

\begin{center}
\begin{tabular}{|l|p{10cm}|}
\hline
\textbf{File} & \textbf{Description} \\
\hline
cache.v & Cache storage structures (data, tag, valid, dirty) \\
clkrst.v & Clock and reset module \\
final\_memory.v & Memory bank implementation \\
four\_bank\_mem.v & Four-banked memory wrapper \\
mem.addr & Example address trace file \\
memc.v & Data storage module used by cache \\
memv.v & Valid-bit storage module used by cache \\
mem\_system\_hier.v & Top-level wrapper used for testing \\
mem\_system\_perfbench.v & Trace-based testbench \\
mem\_system\_randbench.v & Random testbench \\
mem\_system\_ref.v & Reference implementation used by testbench \\
mem\_system.v & Memory system module \textbf{(this is the file you modify)} \\
*.syn.v & Synthesizable memory modules \\
\hline
\end{tabular}
\end{center}

Place the carried-forward cache implementation in:
\begin{quote}
\texttt{./}\textbf{\texttt{demo3}}\texttt{/verilog/phase2\_3/cache\_assoc/}
\end{quote}

Only the associative/final implementation should be copied forward for Phase 3.

\textbf{Your cache controller state diagram should also be carried forward.}

\subsection*{Required student\_custom tests}

Place your custom assembly tests in:

\begin{quote}
\texttt{./assignments/project/testprograms/student\_custom/}
\end{quote}

and add them to:

\begin{quote}
\texttt{./assignments/project/testprograms/student\_custom/all\_phase\_3.list}
\end{quote}

At least two custom tests are required.
Writing more tests is strongly encouraged.

\subsection*{Synthesis}
You should also synthesize your processor and submit the results of synthesis, including the
area, cell, and timing reports. This is similar to what you did for Lab 3.

\subsection*{Submission Contents}

\paragraph{verilog/}

This directory should contain all Verilog and Vcheck files needed to compile and run
your processor. Your processor should be able to compile and run using only the files
present in this directory.

The Phase 3 testing and submission flow is intended to match the rest of the
course workflow: you should run the project test command, then generate your
submission using the standard submission wrapper. Any result collection and
packaging details should be handled automatically by the submission process.

Running the full verification suite may take approximately 40 minutes, so plan accordingly.

\subsection*{Grading and Demonstration}

Submissions will ultimately be graded automatically using the course verification flow for Phase 3.

If failures are detected, students may be asked to meet with the course staff to discuss
partial credit and debugging.

If your processor has known failures, you should bring a short written explanation of
the causes of those failures. Demonstrating an understanding of the problem will help
maximize the credit you receive.

If the complete processor is not functional, you may demonstrate partially working
subsystems such as:

\begin{itemize}
\item Stalling instruction memory alone
\item Stalling data memory alone
\item Stalling instruction and data memory together
\item Direct-mapped instruction cache
\item Direct-mapped data cache
\item Two-way instruction cache
\item Two-way data cache
\end{itemize}

During any demonstration, you should be prepared to explain the behavior of the processor design.

\subsection*{Testing}

Use:
\begin{quote}
\texttt{./run test [-v] project phase\_3 [subtest]}
\end{quote}

\subsection*{Submitting}

Use:
\begin{quote}
\texttt{./run build\_submission project phase\_3}
\end{quote}

\end{document}
